% Annotated reconstruction of the supplied R24, Sections 5--6.
% Uses the fixed valued fields and first-stage statements of D.15--D.16.
\subsection{Arbitrary collar roots and the specified completion}

The next two claims reconstruct the second stage of the supplied
R24 argument \citep[Sections~5--6]{D_OpenAI2026}.
We retain the fields, chosen valuation and embeddings of
D.15--D.16, including $\Omega\subset H\subset E$, the full
$H$-valued point above $\xi$, and the fixed odd prime
$\ell\ne p$. In particular $v=-\log_p|\cdot|$ and
$L_\ell\subset H$ by D.16. The roots below are analytic over
$\Omega$; no descent of those roots to $\bar k$ is required.

\begin{claim}{D.17}{Proved}\label{clm:D17}
Let $m\ge1$, choose an integer
\[
 D>v_p(m)+\frac1{p-1},
\]
and let $\alpha\in\bar k^*$ satisfy $a=v(\alpha)>2D$.
There are a smooth proper connected curve $Q/\bar k$ and a
finite surjective map $Q\to\mathbb P^1_{\bar k}$ whose
horizontal ramification indices divide $\ell$, such that the
coordinate $T$ has an analytic $m$-th root on the entire inverse
image in $Q_\Omega^{\mathrm{an}}$ of
\[
                   A^\circ=\{D<v(T)<a-D\}.
\]
The conclusion holds for every $m$, including powers of $p$
and the case $p=2$.
\end{claim}
\begin{proof}
\emph{Division on the uniformizing torus, on every sheet.}
We first prove the precise period-lattice statement used below,
corresponding to R24, Lemma~5.1.
Suppose $q:T^{\mathrm{an}}\to J^{\mathrm{an}}$ is a toric
uniformization with free discrete kernel $\Lambda\subset T(\Omega)$,
and $q$ is a topological covering and a local analytic isomorphism.
For an analytic map $a_A:A\to T^{\mathrm{an}}$, form
\[
 B=A\times_{q\circ a_A,J^{\mathrm{an}},[m]}J^{\mathrm{an}},
       \qquad \pi:B\to A,\quad b:B\to J^{\mathrm{an}}.
\]
Every character value $\chi\circ a_A\circ\pi$ has an analytic
$m$-th root on all of $B$.
Indeed, locally lift $b$ to $t:B\to T^{\mathrm{an}}$.
The relation $[m]b=q\circ a_A\circ\pi$ gives
\[
                  t^m/(a_A\circ\pi)\in\Lambda.
\]
This is locally constant because the kernel is discrete.
Replacing $t$ by $t\rho$, $\rho\in\Lambda$, changes this
period by $\rho^m$. Thus its class in $\Lambda/\Lambda^m$
is globally defined and gives a finite clopen decomposition
$B=\coprod_{\bar\lambda}B_{\bar\lambda}$.
Choose a representative $\lambda$ of each nonempty class.
On $B_{\bar\lambda}$, a locally constant period correction
makes every local lift satisfy
\[
                         t^m=\lambda(a_A\circ\pi).
\]
Two corrected lifts on an overlap differ by a period $\rho$
with $\rho^m=1$. The free lattice has no torsion, so $\rho=1$.
The corrected lifts therefore glue to one analytic map
$t_{\bar\lambda}:B_{\bar\lambda}\to T^{\mathrm{an}}$.
Choose $d_\lambda\in\Omega^*$ with
$d_\lambda^m=\chi(\lambda)$. Then
$\chi(t_{\bar\lambda})/d_\lambda$ is the required root.
The subsets are disjoint, so these functions give one function
on all of $B$. Neither local connectedness of a chosen sheet
nor a prime-to-$p$ restriction on $m$ has been inserted.

\emph{The auxiliary curve and its split stable model.}
Let $Z$ be the smooth projective normalization of
\[
                         y^\ell=\frac{T(T-1)}{T-\alpha}.
\]
The four branch orders at $0,\alpha,1,\infty$ are
$1,-1,1,-1$. Thus $Z\to\mathbb P^1$ is a connected cyclic
cover of degree $\ell$, totally ramified at those points and
unramified elsewhere. Riemann--Hurwitz gives
\[
                  2g(Z)-2=-2\ell+4(\ell-1),
                     \qquad g(Z)=\ell-1.
\]
All algebraic data descend to a finite extension of $k$.
Over its valuation ring, blow up the line model along
$(T,\alpha)$. The resulting semistable model has two rational
components and one node with chart $TU=\alpha$; the inner
component has coordinate $S=T/\alpha$.
The branch sections are disjoint and smooth. The Kummer
function reduces to $T-1$ on the outer component and
$-S/(S-1)$ on the inner component. Their degree-$\ell$
normalizations are rational, each with exactly two totally
ramified points.

At the node the same function is
\[
                           \frac{T-1}{1-U},
\]
a unit whose reduction is $-1$. Since $\ell$ is invertible,
normalization is \emph{\'etale} there, with $\ell$ split node
charts after a finite constant extension. Away from the node
and the branch sections the unit criterion gives the same
\emph{\'etale} assertion. At a branch section the function or
its inverse is a simple relative parameter, whose Kummer
normalization is smooth. Thus the special fibre consists
exactly of two rational components joined by $\ell$ nodes,
each of thickness $a$. It is stable because $\ell\ge3$;
no regularity of the total space at a node of higher thickness
is required.

The skeleton $\Gamma$ has two vertices and $\ell$ parallel
edges $e_1,\ldots,e_\ell$, all of length $a$.
Its node annuli $A_j$ are disjoint copies of
$\{|\alpha|<|T|<1\}$ with coordinate $T$.
This last all-sheet assertion can also be checked directly:
on that annulus
\[
 \frac{T(T-1)}{T-\alpha}
                =-\frac{1-T}{1-\alpha/T},
\]
and both principal units have convergent $\ell$-th roots.
The full inverse image is therefore $\ell$ copies of the annulus.

\emph{An integral character of slope one.}
Choose the rational branch point over $T=0$ as base point for
an algebraic Abel--Jacobi map $i:Z\to J$.
The rational components imply that the Jacobian is totally
degenerate. Over $\Omega$ its uniformization is
\[
                q:T^{\mathrm{an}}\longrightarrow J_\Omega^{\mathrm{an}},
                  \qquad\ker q=\Lambda.
\]
The uniformization is a topological covering and local analytic
isomorphism with free discrete period lattice, as recalled in
\citet[Sections~4.1--4.2]{D_BakerRabinoff2015}.
That source's Section~4.5.2 identifies the character lattice
with $M=H_1(\Gamma,\Z)$; Proposition~6.1 and Corollary~6.6
identify the tropicalized Abel--Jacobi map with integration on
$\Gamma$ using the edge-length pairing.

An open Berkovich annulus has an interval as skeleton and
open-ball branches; it is locally path connected and simply
connected. Thus $i|_{A_j}$ lifts continuously to the covering
torus, and the local analytic inverses of $q$ make this lift
an analytic map $b_j:A_j\to T^{\mathrm{an}}$.
Orient every edge by increasing $v(T)$. For $j'\ne j$, the
cycle $c_j=e_j-e_{j'}$ is an integral element of $M$.
Let $\chi_j$ be its character and put $u_j=\chi_j\circ b_j$.
Integration of the harmonic form with edge coefficients $c_j$
along a segment of $e_j$ equals the length of that segment.
Consequently
\[
                       \frac{d\,v(u_j)}{d\,v(T)}=1.
\]
This is character evaluation in the stated integral lattice,
not evaluation after inversion of a period matrix. Different
torus lifts add a constant tropical period and do not change
the derivative.

\emph{The middle collar and arbitrary root orders.}
The annular-unit factorization used in D.16 gives
\[
       u_j=c_j'T(1+h_j),\qquad
       h_j=\sum_{n\ne0}h_{j,n}T^n,\qquad c_j'\in\Omega^*.
\]
For clarity, the term $c_j'T$ is the actual degree-one Laurent
coefficient. At a Gauss radius in $|\Omega^*|$, the reductions
of a unit and its inverse are inverse units in the Laurent
polynomial ring, so exactly one monomial dominates.
The exponent cannot change inside the annulus: convergence
on compact subannuli reduces a proposed change to finitely
many terms, and a tie occurs at a radius in $|\Omega^*|$
because $v(\Omega^*)=\Q$. Such a tie is impossible.
The known slope forces the exponent to be one everywhere.

Dominance at all radii $|\alpha|<r<1$, followed by limits
at the endpoints, gives for $n>0$
\[
                         |h_{j,n}|\le1,
                   \qquad |h_{j,-n}|\le|\alpha|^n.
\]
It follows at $s=v(T)$ that
\begin{equation}\label{eq:D17:middle-bound}
 |h_j|\le p^{-\min\{s,a-s\}}<p^{-D}
                      \quad(D<s<a-D).
\end{equation}
Every closed subannulus of this collar has a strict uniform
margin. The logarithm and exponential therefore converge
locally uniformly, and
\[
                b_j'=\exp\bigl(m^{-1}\log(1+h_j)\bigr)
\]
is analytic with $(b_j')^m=1+h_j$.
Indeed $D>1/(p-1)$ gives
$|\log(1+h_j)|\le|h_j|$, and the exponent has valuation
strictly greater than $D-v_p(m)>1/(p-1)$.
These inequalities explicitly include $p=2$ and arbitrary
$p$-power divisors of $m$.

Now form the single algebraic finite \emph{\'etale} cover
\[
                    Z'=Z\times_{i,J,[m]}J\longrightarrow Z.
\]
It is \emph{\'etale} for every $m$ because the ground field
has characteristic zero. On the full inverse image of each
$A_j$, the period-lattice argument, applied to $a_A=b_j$,
gives an $m$-th root of $u_j$. Dividing by the pullback of
$b_j'$ and by an $m$-th root of $c_j'$ gives an $m$-th root
of $T$ over the middle collar. The $A_j$ are disjoint, so the
functions combine on the entire inverse image of $A^\circ$.
Choose any connected component $Q$ of $Z'$ over $\bar k$.
Its map to the connected $Z$ is finite \emph{\'etale} of
positive degree and therefore surjective. It is smooth and
proper, and retains the roots on every inverse-image sheet.
The composite $Q\to\mathbb P^1$ has only horizontal indices
$1$ and $\ell$, because its upper map to $Z$ is \emph{\'etale}.
All covers used are algebraic over $\bar k$; their analytic
roots and uniformization lifts need only exist over $\Omega$.
\end{proof}

\begin{claim}{D.18}{Proved}\label{clm:D18}
For the fixed prolongation of the real-valued rank-one norm in
D.15, the specified subfield $K^{\mathrm{sep}}$ of its completion
$E$ is contained in $H$, the closure of $\widetilde K$ inside
$E$. In particular $H=E$.
\end{claim}
\begin{proof}
We reconstruct R24, Section~6, with the field copies retained
throughout \citep{D_OpenAI2026}.
Let $F/K_0$ be any finite Galois extension inside the fixed
$K^{\mathrm{sep}}$, and let $C\to X_{\bar k}$ be its smooth
proper curve cover. Its branch set $B$ is finite. If it is
empty, then $F\subset\widetilde K\subset H$ by definition.
Suppose $B\ne\varnothing$.

For every $b\in B$, choose a rational local parameter $t_b$
and a disk $B_b$ as in D.15: the disks are pairwise disjoint,
avoid $\xi$, and the full cover above each is a disjoint union
of maps $t_b=z_c^{e_c}$ on disks. Put
$m_b=\operatorname{lcm}_{c\mid b}e_c$.
Choose $\beta_b\in\bar k^*$ whose absolute value is smaller
than the radius of $B_b$ in the coordinate $t_b$.
Choose an integer
$D_b>v_p(m_b)+1/(p-1)$ and
$\alpha_b\in\bar k^*$ with $a_b=v(\alpha_b)>2D_b$.
D.17 gives a cover of the $T$-line with horizontal indices
dividing $\ell$ and a full inverse-image root of $T$ on
$D_b<v(T)<a_b-D_b$.

Pull this cover back by the finite nonconstant morphism
$X_{\bar k}\to\mathbb P^1_{\bar k}$ defined by
$T=t_b/\beta_b$, normalize, and take a connected component
$W_b\to X_{\bar k}$.
The horizontal-index statement of D.15 shows that its indices
still divide $\ell$. The analytic roots pull back to every
point of this normalized component over the collar
\[
 |\beta_b|p^{-(a_b-D_b)}<|t_b|<|\beta_b|p^{-D_b}
                         \quad\text{inside }B_b.
\]
Multiplication by an $m_b$-th root of $\beta_b$ in $\Omega$
gives a root of $t_b$ there.

Choose embeddings of the finitely many fields $\bar k(W_b)$
into the fixed $K^{\mathrm{sep}}$ over $K_0$, and form their
compositum. Let $W\to X_{\bar k}$ be its smooth proper
connected curve. It dominates every $W_b$, so all the roots
persist on the full inverse collars in $W_\Omega^{\mathrm{an}}$.
D.15's compositum calculation keeps every horizontal index a
divisor of $\ell$; it uses the residue-characteristic-zero
closed-point valuations, not the vertical reduction valuations.

Apply the all-sheet patching construction of D.15 with
$Y=W_\Omega$ and the just-constructed collars. It gives a
connected finite \emph{\'etale} cover
$V_\Omega\to W_\Omega$ and an analytic map
\[
                  (V_\Omega^{\mathrm{an}})|_U
                          \longrightarrow C_\Omega^{\mathrm{an}}
                         \quad\text{over }X_\Omega^{\mathrm{an}},
\]
where $U$ is the exterior of the smaller closed branch disks
and $\xi\in U$.
Since $W$ is proper and $\bar k\subset\Omega$ are algebraically
closed fields, invariance of finite \emph{\'etale} covers under
this extension gives a connected finite \emph{\'etale} cover
$V\to W$ and an isomorphism
\[
                         V\times_{\bar k}\Omega\simeq V_\Omega
                                  \quad\text{over }W_\Omega
\]
\citep[Expos\'e~X, Corollary~1.8]{E_SGA1}.
Only this finite \emph{\'etale} cover is descended, not the
analytic exterior map or its root trivializations.

The composite $V\to X_{\bar k}$ has horizontal indices
still dividing $\ell$, because its upper map is \emph{\'etale}.
Extend the selected embedding of $\bar k(W)$ to
$\bar k(V)$ inside $K^{\mathrm{sep}}$.
Its Galois closure has the same index bound by D.15, so the
image of $\bar k(V)$ lies in $L_\ell$ and hence in $H$ by D.16.
This membership does not require $V\to X$ itself to be Galois.

The evaluation statement of D.15 now gives a $K_0$-embedding
$F=\bar k(C)\hookrightarrow H$. Explicitly, the chosen
embedding of $\bar k(V)$, together with $\Omega\subset H$,
gives an $H$-valued point of $V_\Omega$ over the originally
specified $H$-valued point of $X_\Omega$. Its analytic image
lies over $\xi\in U$. Evaluate the exterior analytic map
using the induced isometric map of completed residue fields.
The resulting point of $C_\Omega(H)$ projects to the generic
point of $X_{\bar k}$ with the specified inclusion $K_0\to H$.
Since $C\to X_{\bar k}$ is finite, its support on $C$ is
generic as well. It therefore gives the stated field embedding.
Keeping the entire $H$-valued point makes the argument valid
even if $\xi$ itself is an $\Omega$-rational analytic point.

Finally this embedding lands in the specified copy of $F$.
For each $x\in F$, its minimal polynomial over $K_0$ splits
in $F$ because $F/K_0$ is normal. Its image in $H\subset E$
is one of these same roots: a split polynomial over a field
cannot acquire additional roots in a larger field. Thus the
image is contained in $F$, and equality follows from
$[F:K_0]$. Consequently the chosen $F\subset E$ is contained
in $H$.

Every element of $K^{\mathrm{sep}}$ lies in such a finite
Galois extension of $K_0$, since $K^{\mathrm{sep}}$ is an
algebraic closure of $K_0$ in characteristic zero. Taking their
union proves $K^{\mathrm{sep}}\subset H$ inside $E$.
The subfield $H$ is closed in $E$ and contains its dense subfield
$K^{\mathrm{sep}}$, so $H=E$.
\end{proof}
